\chapter{Conclusion and Future Work}

\section{Conclusion}

This dissertation set out to investigate whether a domain-fine-tuned
large language model, embedded inside an agentic self-correction loop,
can autonomously produce syntactically valid and semantically faithful
end-to-end user-interface test scripts directly from natural-language
user stories. The motivation was the persistent quality-assurance
bottleneck in modern software teams, where manual authoring and
maintenance of end-to-end test scripts consumes disproportionate
engineering effort relative to the value delivered. Eight phases of work
were undertaken to answer this question: a curated, framework-balanced
dataset of 279 Jira-style user stories was constructed and validated;
two commercial large language models were evaluated as zero-shot
reference systems; two open-weight models, Phi-3 Mini and Gemma 4 E4B,
were adapted through QLoRA parameter-efficient fine-tuning into
framework-isolated adapters for Cypress and Playwright; a
Build--Measure--Assess--Decide agentic correction loop was implemented
to iteratively refine generated scripts against a composite quality
threshold; and a comprehensive statistical evaluation was carried out
comparing all eight resulting systems on the complete held-out dataset.

The evaluation reported in Chapter 5 answers the research question in
the affirmative, with three specific, statistically substantiated
conclusions. First, fine-tuning produces a large and highly consistent
improvement in ground-truth semantic alignment over zero-shot prompting:
the ranking Gemma 4 E4B \textgreater{} Phi-3 Mini \textgreater{}
commercial baselines holds without a single exception across both
frameworks and across all sixteen functional categories in the dataset,
and every one of the ten paired Wilcoxon signed-rank comparisons
conducted was significant at p \textless{} $10^{-22}$. Second, this
improvement is not obtained at the expense of syntactic reliability:
once evaluated on the complete dataset rather than a partial sample,
fine-tuned and baseline systems are statistically indistinguishable on
syntax validity ($\chi^2$ = 0.00, p = 1.00). Third, the agentic correction
loop itself is an effective mechanism for closing the residual gap
between a first-attempt generation and an acceptable script, with the
strongest configuration --- Gemma 4 E4B on Cypress --- reaching a 100\%
acceptance rate across all 279 held-out stories.

Beyond the headline quantitative result, this work makes three further
contributions that are methodological rather than purely empirical. The
evaluation design pairs every one of the 279 stories across all eight
systems, enabling item-level non-parametric significance testing rather
than comparison of aggregate means alone --- a stronger design than is
typical in comparable code-generation studies, where baseline and
treatment systems are often evaluated on non-identical or non-paired
samples. Second, two measurement artefacts were identified and corrected
during the course of the evaluation itself: an unstripped Markdown code
fence that had suppressed Claude Haiku's apparent Playwright
syntax-validity score by more than thirty percentage points, and an
incomplete-dataset sample that had made fine-tuned syntax reliability
appear to lag baseline reliability by a statistically significant
margin, a gap that closed entirely once the full 279-story dataset was
scored. Both corrections are reported transparently in Chapter 5 rather
than silently absorbed into the final figures, illustrating a broader
methodological point for automated code-generation research: an
unexpected per-system result should be treated as a candidate
measurement artefact and investigated before it is accepted as a genuine
finding. Third, the category-level analysis in Section 5.7 demonstrates
that a lexical-overlap metric such as ROUGE-L can systematically
understate model quality in categories whose reference scripts rely on
idiosyncratic, non-derivable naming conventions, a qualification that
should inform how future work in this area interprets ROUGE-L-based
comparisons.

\section{Limitations}

Several limitations qualify the scope of the conclusions reported above
and should be considered when generalising the findings beyond the
conditions under which this study was conducted.

\begin{itemize}
\item
  Dataset scale and provenance. The 279-story dataset, while balanced
  across sixteen functional categories and three complexity levels, is a
  single synthetically constructed corpus rather than a sample of user
  stories drawn from an operating production backlog. Although its scale
  is consistent with established parameter-efficient fine-tuning
  practice for LoRA-style adapters, a larger and organisationally
  diverse dataset would strengthen claims of generalisability to
  arbitrary real-world Jira tickets.
\item
  Framework coverage. Evaluation is restricted to Cypress and
  Playwright, the two most widely adopted browser-based end-to-end
  testing frameworks at the time of this study. The findings do not
  directly extend to other test automation paradigms such as Selenium
  WebDriver, native mobile frameworks such as Appium, or API-level
  contract testing, each of which carries a different syntactic and
  semantic profile.
\item
  Model scale. Both fine-tuned base models --- Phi-3 Mini and Gemma 4
  E4B --- were selected in part for their feasibility on Apple Silicon
  hardware operating under a 48 GB unified-memory budget. Larger
  open-weight or proprietary models may exhibit different fine-tuning
  dynamics and zero-shot baseline quality; the results reported here
  should be read as evidence for models in this parameter range, not as
  an upper bound on what fine-tuning can achieve at larger scale.
\item
  Single fine-tuning run per configuration. Each of the four adapters
  was trained once to convergence rather than across multiple random
  seeds. Training-run-to-training-run variance was therefore not
  quantified, and the reported accept rates and ROUGE-L scores reflect
  one instantiation of each adapter rather than a distribution over
  repeated training runs.
\item
  Token-budget-bound rejections. A minority of Phi-3 Mini rejections,
  most visible on Playwright, were attributable to the generation token
  ceiling being insufficient for unusually long reference scripts even
  after the correction loop's automatic escalation. This is a
  configuration limit of the present loop rather than evidence of a
  structural incapacity in the model, but it does mean the reported
  accept rates slightly understate what a more generous or adaptively
  sized token budget could achieve.
\item
  Metric-based rather than execution-based validation. Quality was
  assessed through static syntax checking and lexical similarity to a
  reference script; no generated script was executed against a live
  application under test. A script can therefore be judged syntactically
  valid and lexically close to the reference while still failing at
  runtime --- for example, if a selector it references does not exist in
  the target application's actual DOM. This is discussed further as a
  priority direction for future work below.
\end{itemize}

\section{Future Work}

The limitations identified above translate directly into concrete
directions for extending this research.

\subsection{Execution-Based Validation}

The most consequential extension is to move beyond static and lexical
evaluation toward execution-based validation: running each generated
script against a live or containerised instance of the target
application and recording pass, fail, or runtime-error outcomes. This
would close the remaining gap between ``the script looks correct'' and
``the script actually verifies the intended behaviour,'' and would allow
the BMAD loop's Measure stage to incorporate a genuine
functional-correctness signal alongside the current static heuristics. A
natural implementation path is a sandboxed headless-browser execution
stage inserted between the current syntax check and the composite-score
computation, feeding pass/fail outcomes back as an additional term in
the acceptance threshold.

\subsection{Human Evaluation}

A structured human evaluation, in which practising QA engineers rate a
sample of generated scripts for readability, maintainability, and
fitness for direct use in a real codebase, would complement the
automated metrics used in this dissertation. ROUGE-L and token-F1
capture surface similarity to a reference implementation but cannot
assess whether an engineer would accept a generated script into a pull
request with minimal edits --- a criterion that ultimately determines
the practical value of the system studied here.

\subsection{Broader Framework and Domain Coverage}

Extending the framework-isolated fine-tuning methodology used here to
Selenium WebDriver, Appium-based mobile test automation, and API-level
contract testing frameworks such as Postman or REST Assured would test
whether the observed fine-tuning advantage generalises beyond
browser-based UI testing, and would allow a single organisation's full
test-automation surface to be addressed by variants of the same agentic
loop.

\subsection{Larger-Scale and Multi-Seed Training}

Repeating the QLoRA fine-tuning procedure across multiple random seeds
per adapter would quantify training-run variance and strengthen the
statistical claims made in Chapter 5. Scaling the dataset beyond 279
stories per framework, and sourcing a portion of it from real,
anonymised production backlogs rather than synthetic generation, would
additionally test whether the fine-tuning gains reported here persist on
more heterogeneous and less internally consistent user-story phrasing.

\subsection{Adaptive Token Budgeting}

The token-budget-bound rejections identified in Section 6.2 could be
addressed directly by predicting an appropriate generation length from
the user story's category and complexity label before the first
generation attempt, rather than escalating the token ceiling reactively
only after a truncated first attempt --- a refinement that would likely
narrow the residual accept-rate gap between Phi-3 Mini and Gemma 4 E4B
on Playwright.

\subsection{Multi-Agent Specialisation}

The current architecture applies a single agentic loop uniformly across
all categories of user story. A natural extension is to decompose the
loop into specialised sub-agents --- for instance, a dedicated agent for
accessibility-focused assertions, or one specialised in
responsive-design viewport handling, the category identified in Chapter
5 as the most metric-sensitive --- coordinated by a routing mechanism
that dispatches each incoming story to the most appropriate specialist
before the shared Build--Measure--Assess--Decide cycle proceeds.

\section{Concluding Remarks}

Taken as a whole, this dissertation provides statistically robust
evidence that parameter-efficient fine-tuning, combined with an agentic
self-correction loop, can substantially improve an open-weight language
model's ability to translate natural-language user stories into usable
end-to-end test scripts, without sacrificing syntactic reliability
relative to considerably larger zero-shot commercial systems. The scale
of the improvement, its consistency across every category and framework
tested, and the rigour with which measurement artefacts were identified
and corrected during the evaluation process together support the
conclusion that this is a genuine and practically meaningful effect
rather than an outcome of favourable aggregation. The limitations and
future directions identified in this chapter define a clear path by
which the present findings could be extended toward a system suitable
for evaluation in an operating software delivery pipeline.
